\section{議論}\label{sec:discussion}

\subsection{未解決の問題}
\begin{enumerate}
\item \textbf{凍結した状態の多重項}：多箱の $c = 1$ 部分は $\sigma$ 対称でなく、$G$ では展開できない。$G = T_+T_-$ の因子 $T_\pm$ を使っても、$\sigma$ 対称でない部分は展開できるようにならない（§\ref{sec:index} の注意）。
向き付けられない面の refined の取っ手の因子を幾何から決める必要がある。
\item \textbf{$E$ と $C$ の分離}：指数からは決まらない（命題 \ref{prop:EC}）。Hopf の $(4,4)$ と結び目の 2 箱で未決定である。
固定セクター（被覆の環状の弦、商では Möbius の帯）の Hilbert 空間への $\tilde I$ の作用を、指数ではなく状態の空間の上で構成する必要がある。
\item \textbf{普遍因子の意味}：$(\ba^2 - \bq^2)$ は SO(2) の点で、$(\bt^2 - 1)$ は $t = q$ と偶奇で説明できる。ただし、$\nu \to \tnu$ のずれの向きが一意に決まる理由と、閉弦の Kähler 変数の三重化との関係は未解明である。
\item \textbf{他の結び目}：非トーラス結び目では、refined SO の値に定義が無い（refined CS の $S, T$ が使えない）。
固定 $N$ では、SO(4) $=$ SU(2)$^2$、SO(6) $=$ SU(4) の同一視で A 型の入力（KNTZ \cite{KNTZ19} の長方形彩色の式など）を使える。二重ツイスト結び目 7 個で、$N = 2$ の SO 側（$|R| \le 3$）の hatted の整数性と一様な規格化が成り立つことを確かめた。
大きな $N$ の値には、D 型の DAHA の iterated torus 結び目への拡張が要る。
\item \textbf{Sp と B 型}：本論文は SO$(2N)$（D 型）に限る。閉弦の B 型（AM）は $c = 1$ だけで、分母に $(t^{n/2} + t^{-n/2})$ をもつ。命題 \ref{prop:alld} の分類で B 型をどう読むかは調べていない。
\end{enumerate}

\subsection{KN との関係}
被覆（向き付けられた部分）の brane--brane のセクターは、KN の $G$ 展開と正値性をそのまま満たす。
破れは、$\tilde I$ に関わるセクター（brane--image とクロスキャップ）に限られる。
したがって、refined の Kauffman 不変量の正値性の問題は、KN の枠組みに orientifold の偶奇の構造を加える問題に帰着する（予想）。
